%%%BLOCK id=001-02 ch=11 sec=洛伦兹力分解·微元法 kind=derivation alt=none cont=next y=0.15-0.23
分解洛伦兹力：
\begin{align*}
qv_xB &= ma_y\\
qBx &= m\Delta v_y % CHECK: 原稿"qBx=m△v_y"，△v_y 前的小三角符号较小；左边的x疑为位移(Δx)，照原样转录
\end{align*}
%%%END
%%%BLOCK id=014-07 ch=11 sec=地磁场 kind=knowledge alt=none cont=none y=0.82-0.96
\textbf{磁场}

%FIG p014_d 0.11 0.828 0.285 0.958
\notefig[0.25]{figures/p014_d.png}
南极\ $B$向上

北极\ $B$向下

赤道\ 平行于地面向北

南半球：$B_1$向北、$B_2$向上

北半球：$B_1$向北、$B_2$向下
%%%END
%%%BLOCK id=015-01 ch=11 sec=安培力·有效长度 kind=formula alt=none cont=none y=0.04-0.19
\textbf{安培力}

$F_A=BIL\sin\theta$

%FIG p015_a 0.32 0.03 0.66 0.184
\notefig[0.5]{figures/p015_a.png}
% 图中：弧线A→B；直线A→B；闭合圆环；与电容器相连的圆环（内有“d”和一条虚线）。以下为各图旁的标注：
$L_{\text{有效}}=\overline{AB}$\qquad $L_{\text{有}}=\overline{AB}$\qquad $L=0$\qquad $L_{\mbox{\unclear{上}}}=d$\quad $L_{\mbox{\unclear{下}}}=d$

若$B$、$I$不垂直\quad 分解$B$
%%%END
%%%BLOCK id=015-02 ch=11 sec=安培力·微观表达 kind=derivation alt=none cont=none y=0.19-0.27
\[
F_A=BIL=BL\,nqSv=LSn\cdot Bvq=\overset{\text{电子总数}}{\underline{LSn}}\cdot F_{\text{洛}}
\]
$n$ 单位体积电子数
%%%END
%%%BLOCK id=015-03 ch=11 sec=电流的磁效应·奥斯特实验 kind=knowledge alt=none cont=none y=0.27-0.33
奥斯特实验：电流产生磁场

电流应水平南北放置，小磁针应放到电流正上/下方
%%%END
%%%BLOCK id=015-04 ch=11 sec=电流的磁场·直线电流 kind=formula alt=none cont=none y=0.33-0.37
\[
B=K\frac{I}{r}=\frac{\mu_0 I}{2\pi r}
\]
$\mu_0$ 真空磁导率\quad $\dfrac{\mathrm{m\cdot T}}{\mathrm{A}}$
%%%END
%%%BLOCK id=015-05 ch=11 sec=安培力·通电导线间作用力 kind=method alt=none cont=none y=0.355-0.455
通电导线间作用力\quad 同吸异斥

先判磁场后判力
%FIG p015_b 0.463 0.358 0.725 0.452
\notefig[0.4]{figures/p015_b.png}
%%%END
%%%BLOCK id=015-06 ch=11 sec=安培力·做功与冲量 kind=formula alt=none cont=none y=0.455-0.57
\textbf{安培力做功}：
\[
\left\{\begin{aligned}
&\text{恒力}\quad \pm BIL\cdot x\\
&\text{变力}\quad \pm W_{\text{安}}
\end{aligned}\right.
\]
\textbf{安培力冲量}：
\[
\left\{\begin{aligned}
&\text{恒力}\quad \pm BILt=\pm BLq\\
&\text{变力}\quad \pm I_{A}=\pm BLq
\end{aligned}\right.
\]
%%%END
%%%BLOCK id=015-07 ch=11 sec=磁偏转·轨迹圆心半径 kind=method alt=none cont=none y=0.565-0.66
\textbf{磁偏转}：
\begin{itemize}
\item 画轨迹
\item 找圆心：
\begin{itemize}
\item 速度作垂线过圆心
\item 弦作中垂线过圆心
\item 速度偏转角补角分线过圆心
\end{itemize}
\item 求半径
\end{itemize}
%%%END
%%%BLOCK id=015-08 ch=11 sec=动态圆 kind=summary alt=none cont=none y=0.66-0.925
% 手写大括号括住下面四行（静态圆～汇聚圆/发散圆）；“汇聚圆/发散圆”下有一条带箭头的下划线
\begin{center}
\begin{tabular}{lcc}
 & $v_0$方向 & $r$\\
静态圆 & 不变 & 不变\\
放缩圆 & 不变 & 变\\
旋转圆 & 变 & 不变\\
汇聚圆/发散圆 & \multicolumn{2}{l}{$\to$ $r_{\text{粒子}}=R_{\text{磁场}}$}
\end{tabular}
\end{center}
%FIG p015_c 0.05 0.735 0.72 0.918
\notefig[0.75]{figures/p015_c.png}
\textbf{图中标注}：②、①，$v_0$，$r=\dfrac{mv}{qB}$，圆心圆，远点圆，$d=2r=\dfrac{2mv}{qB}$
%%%END
%%%BLOCK id=015-09 ch=11 sec=动态圆 kind=summary alt=none cont=next y=0.92-0.95
\textbf{平移圆}
%%%END
%%%BLOCK id=016-01 ch=11 sec=平移圆 kind=method alt=none cont=none y=0.10-0.20
\textbf{平移圆}\quad $V_0$定\quad 入射位置改变
%FIG p016_a 0.10 0.12 0.36 0.195
\notefig[0.3]{figures/p016_a.png}
%%%END
%%%BLOCK id=016-02 ch=11 sec=时间极值 kind=method alt=none cont=none y=0.19-0.30
\textbf{时间极值}\quad $t=\dfrac{\theta m}{qB}=\dfrac{s}{V_0}$\quad 时间最长\ $\theta$最大\ 圆心角最大
% CHECK: 分子手写潦草，像 θm（也可能是 2πm）；"时间最长"后面一字同样像 θ（也可能是 弧）
\begin{itemize}
\item $V_0$定值\quad 最长弦/弧（劣）
\item $V_0$变化\quad 最大圆心角
\end{itemize}
% 原稿：以上两行左侧有大括号；"最长弦/弧（劣）"下方有一条波浪线向右延伸至"最大圆心角"上方
%%%END
%%%BLOCK id=016-03 ch=11 sec=单边有界入射 kind=method alt=none cont=none y=0.30-0.42
\textbf{单边有界入射}
%FIG p016_b 0.05 0.335 0.205 0.415
\notefig[0.25]{figures/p016_b.png}
角入角出\quad 速偏角$=$圆心角$=2$倍弦切角
%FIG p016_c 0.49 0.315 0.65 0.390
\notefig[0.25]{figures/p016_c.png}
\red{圆心圆}\quad \red{定速不定向问题}
% 原稿：红色弧线从"圆心圆"（浅粉底）右上角向下弯至"弦切角"右侧，连接这两处
%FIG p016_d 0.675 0.315 0.80 0.381
\notefig[0.25]{figures/p016_d.png}
垂直入射即半圆
%%%END
%%%BLOCK id=016-04 ch=11 sec=单边有界出射型 kind=method alt=none cont=none y=0.42-0.60
\textbf{单边有界出射型}\quad \red{O离板距$d$}\quad \red{光屏}
%FIG p016_e 0.22 0.415 0.52 0.52
\notefig[0.35]{figures/p016_e.png}
到不了\unclear{边界}交点，\red{到切点，已撞板了}
% CHECK: "到不了"与"交点"之间被光屏横线盖住的字看不清，像 边界/切
% 原稿：图右侧有一个竖直大括号，旁注 3m
$3\unit{m}$\quad $r=2.5\unit{m}$
%FIG p016_f 0.825 0.42 0.96 0.48
\notefig[0.25]{figures/p016_f.png}
% 红色小图（直角三角形）文字：斜边 R，竖直边 d-R（符号像 + 或 -，按物理意义取 -），水平边 sqrt(R^2-(d-R)^2)
% CHECK: 红字 d 与 R 之间的符号手写像 + 也像 -
\begin{itemize}
\item 最值相切：最远点为轨迹与边界切点
\item 最值相交：最远点为直径和边界交点
\item 能交则交，不能交则切。
\end{itemize}
% 原稿：以上三行左侧有大括号
%%%END
%%%BLOCK id=016-05 ch=11 sec=圆形磁场·汇聚发散 kind=method alt=none cont=none y=0.60-0.80
\textbf{汇聚\ 发散}
%FIG p016_g 0.02 0.64 0.52 0.815
\notefig[0.5]{figures/p016_g.png}
$V$大小变化\ 方向定

$B_1=B_2$\ 点入平出

平进点出
%%%END
%%%BLOCK id=016-06 ch=11 sec=径入径出 kind=method alt=none cont=none y=0.80-0.95
%FIG p016_h 0.04 0.795 0.27 0.97
\notefig[0.3]{figures/p016_h.png}
径入径出
%FIG p016_i 0.49 0.82 0.745 0.955
\notefig[0.3]{figures/p016_i.png}
化单边处理
%%%END
%%%BLOCK id=017-01 ch=11 sec=环形磁场 kind=method alt=none cont=none y=0.06-0.29
\textbf{环形磁场}
%FIG p017_a 0.10 0.095 0.37 0.29
\notefig[0.3]{figures/p017_a.png}
最值相切\quad 径入径出
\[ (R-R_2)^2=R_2^2+r^2 \qquad R_2=\frac{mV_2}{qB}\ \Rightarrow\ V_2 \]
\[ (R_1+r)^2=R^2+R_1^2 \qquad R_1=\frac{mV_1}{qB}\ \Rightarrow\ V_1 \]
%%%END
%%%BLOCK id=017-02 ch=11 sec=圆形磁场·汇聚发散 kind=method alt=none cont=none y=0.29-0.60
\textbf{圆形磁场汇聚发散}\quad $R_{\text{轨迹}}=R_{\text{磁场}}$

点入平出\quad 平入点出
%FIG p017_b 0.095 0.368 0.352 0.595
\notefig[0.4]{figures/p017_b.png}
%FIG p017_c 0.41 0.362 0.67 0.545
\notefig[0.35]{figures/p017_c.png}
菱形 $\Rightarrow BO_2\parallel AO_1\Rightarrow V_{\text{末}}\parallel L$
%%%END
%%%BLOCK id=017-03 ch=11 sec=类圆形磁场 kind=method alt=none cont=none y=0.59-0.80
\textbf{类圆形磁场}（半圆、\unclear{$\frac14$圆}、扇形、树叶形）

粒子没有经过的磁场，是可有可无的磁场
%FIG p017_d 0.09 0.665 0.30 0.80
\notefig[0.25]{figures/p017_d.png}
入射范围角$=$圆心角
\[ S_{\mathrm{tot}}=2\left(S_{\text{扇}}-S_{\triangle}\right) \]
% CHECK: 下标手写为"ot"（可能是 tot 或 总）；S 的第二个下标像 角/扇，按物理意义取 扇
%%%END
%%%BLOCK id=017-04 ch=11 sec=双直线边界 kind=method alt=none cont=next y=0.80-1.0
\textbf{双直线边界}

最值相切\quad 能从上口射出
%FIG p017_e 0.04 0.884 0.208 1.0
\notefig[0.25]{figures/p017_e.png}
定向不定速

$\checkmark$最\unclear{小}（最\unclear{小}$V$）\quad 圆的放缩
% CHECK: 第一个对勾后的"最?"与括号内"最?V"手写潦草，按物理意义（相切对应最小速度）取 小
能从下口射出的

$\checkmark$最大（最\unclear{大}$V$）
% CHECK: 括号内"最?V"手写潦草，像 木（大/小 均有可能），按物理意义取 大
%%%END
%%%BLOCK id=018-01 ch=11 sec=双直线边界·最值相切 kind=derivation alt=none cont=prev y=0.02-0.37
$AD=2R$
%FIG p018_a 0.115 0.045 0.32 0.28
\notefig[0.3]{figures/p018_a.png}
\textbf{最值相切}
%FIG p018_b 0.365 0.10 0.45 0.168
\notefig[0.25]{figures/p018_b.png}
\[ SB^2=R^2-(d-R)^2 \]
\[ SC^2=OT^2=R^2-(d-R)^2 \]
\[ R<d<2R \]
%FIG p018_c 0.68 0.035 0.89 0.285
\notefig[0.3]{figures/p018_c.png}
\[ SB^2=R^2-(R-d)^2 \]
\[ SC^2=R^2-(R-d)^2 \qquad \because SC=SB \]
故\ $SD=2SC$
% 原稿：最后一行"故 SD=2SC"下方有红色下划线
%%%END
%%%BLOCK id=018-02 ch=11 sec=双直线边界·时间 kind=method alt=none cont=none y=0.37-0.60
\textbf{双直线边界}\ $t$\quad 从右边界射出\quad $V$\unclear{同}$\to$弦长最短$\to$圆心角最小
%FIG p018_d 0.12 0.455 0.222 0.595
\notefig[0.25]{figures/p018_d.png}
\[ \sin\theta=\frac{\frac12 d}{R} \]
\[ t_{\min}=\frac{2\theta m}{qB} \]
% 原稿：图右侧另有两条空的竖直长线（未画完的双直线边界图），无其它内容
%%%END
%%%BLOCK id=018-03 ch=11 sec=反向磁场 kind=method alt=none cont=none y=0.60-0.80
\textbf{反向磁场}
%FIG p018_e 0.09 0.637 0.30 0.782
\notefig[0.3]{figures/p018_e.png}
若\ $B_1=B_2$\ $L_1=L_2$\ \underline{则初末速度平行}
% 原稿：此行"则初末速度平行"下方为红色下划线

若\ $B_1\ne B_2$\ $L_1\ne L_2$\ 若初末速度平行\ 则\ \underline{$B_1L_1=B_2L_2$}
% 原稿：此行"B1L1=B2L2"下方为红色下划线，红线向右延伸至页边

速偏角$\alpha_1$，速偏角$\alpha_2$\quad $\alpha_1=\alpha_2\ \overset{\text{相似}}{\Longleftrightarrow}\ \dfrac{R_1}{L_1}=\dfrac{R_2}{L_2}$\quad $R=\dfrac{mv}{qB}$
%%%END
%%%BLOCK id=018-04 ch=11 sec=隐形磁场 kind=method alt=none cont=none y=0.80-0.94
\textbf{隐形磁场}
%FIG p018_f 0.17 0.822 0.42 0.942
\notefig[0.3]{figures/p018_f.png}
\begin{itemize}
\item 最小面积：圆\quad 以入射出射弦为直径的圆
\item 三角形：等边$\triangle$\ 等边三角形
\item 矩形：入射出射弦为长，切弧成形
\item 不定形：粒子未经的地方去除
\end{itemize}
% 原稿：以上四行左侧有大括号；"矩形"一行右端有一个小图（矩形内含一段圆弧）
%FIG p018_g 0.925 0.87 0.995 0.905
\notefig[0.25]{figures/p018_g.png}
%%%END
%%%BLOCK id=018-05 ch=11 sec=找圆心 kind=method alt=none cont=none y=0.94-1.0
\textbf{找圆心}
\begin{itemize}
\item 速度在磁场内：速度作垂线找圆心
\item 速度不在磁场内：速偏角的补角平分线过圆心
\end{itemize}
% 原稿：以上两行左侧有大括号
%%%END
%%%BLOCK id=019-01 ch=11 sec=多解·直线边界吃弦长 kind=method alt=none cont=none y=0.03-0.27
\fbox{多解}\ \textbf{直线吃弦长}

\red{反向磁场}
%FIG p019_a 0.065 0.084 0.38 0.178
\notefig[0.4]{figures/p019_a.png}
\[ (2R_1+2R_2)\,n=L \]
\[ \text{或}\ 2R_1(n+1)+2R_2\,n=L \]

\red{同向磁场}
%FIG p019_b 0.405 0.067 0.66 0.168
\notefig[0.35]{figures/p019_b.png}
\[ (2R_2-2R_1)\,n=L \]
\[ \text{或}\ (n+1)\,2R_2-2R_1\cdot n=L \]
%%%END
%%%BLOCK id=019-02 ch=11 sec=多解·圆弧吃角度 kind=method alt=none cont=none y=0.27-0.50
\textbf{圆弧吃角度}\ \fbox{多解}
%FIG p019_c 0.13 0.314 0.33 0.44
\notefig[0.3]{figures/p019_c.png}
要回到出发点..
\[ \theta\cdot n=2\pi \]
\[ \begin{cases}\theta=\dfrac{2\pi}{n}\quad(n=2,3,\dots)\\[1ex] r=R\cdot\tan\dfrac{\theta}{2}\end{cases} \]
\[ \therefore\ r=R\tan\frac{\pi}{n}\qquad r=\frac{mv}{qB}\qquad \therefore\ V=\frac{qBR\tan\frac{\pi}{n}}{m} \]
%%%END
%%%BLOCK id=019-03 ch=11 sec=复合场·电磁组合场 kind=summary alt=none cont=none y=0.50-0.73
\textbf{电磁组合场}\quad $qE$\ $qvB$
\begin{itemize}
\item 静止\quad $V_0=0$\quad $a=0$\quad $qV_0B=0$\quad $qE\ne0$\quad 不会静止
\item 直线
  \begin{itemize}
  \item 匀速\quad 必为匀速\quad $qvB=qE$\quad $V=\dfrac{E}{B}$
  \item 变速
    \begin{itemize}
    \item $a$定\quad $\times$
    \item $a$变\quad $\times$
    \end{itemize}
  \end{itemize}
\item 曲线
  \begin{itemize}
  \item 类平抛\quad $\times$
  \item 圆周\quad $\times$
  \item 摆线\quad 螺旋线
  \end{itemize}
\end{itemize}
% 原稿：静止/直线/曲线三项左侧有大括号；"直线"下的匀速/变速、变速下的a定/a变、曲线下的三项各有小括号；
% 圆周右侧的 × 位置略偏下，介于"圆周"与"摆线"两行之间，这里归给圆周
%%%END
%%%BLOCK id=019-04 ch=11 sec=复合场·三场叠合 kind=summary alt=none cont=next y=0.73-0.95
\textbf{三场叠合}\quad $mg$\ $qE$\ $qvB$
\begin{itemize}
\item 静止\quad $V_0=0$\quad $qV_0B=0$\quad $mg=qE$
\item 匀圆（纯磁场）\quad 电场力\unclear{重}力抵消\quad $qE=mg$
\item 直线\quad 匀直
\item 螺旋线\quad 配速法
\end{itemize}
% 原稿：以上四项左侧有大括号
%%%END
%%%BLOCK id=020-01 ch=11 sec=复合场·配速法 kind=method alt=none cont=prev y=0.04-0.31
\textbf{曲线}\quad 恒力$+qvB$
\begin{itemize}
\item 洛伦兹力不做功$\to$动能定理
\item 配速法\quad （垂直于恒力方向分解速度$\to$抵消恒力）
\end{itemize}
% 原稿：以上两行左侧有大括号
%FIG p020_a 0.07 0.16 0.26 0.265
\notefig[0.25]{figures/p020_a.png}
$\vec{V_1}+\vec{V_2}=V_0$
%FIG p020_b 0.43 0.18 0.63 0.265
\notefig[0.25]{figures/p020_b.png}
\[ V_1=\frac{mg}{qB}\qquad qV_1B=mg \]
%FIG p020_c 0.64 0.18 0.75 0.305
\notefig[0.25]{figures/p020_c.png}
\[ V_2=\sqrt{V_0^2+\left(\frac{mg}{qB}\right)^2} \]
\[ V_{\min}=V_2-V_1 \]
\[ V_{\max}=V_2+V_1 \]
%%%END
%%%BLOCK id=020-02 ch=11 sec=复合场·速度选择器 kind=method alt=none cont=none y=0.31-0.44
\textbf{速度选择器}\quad $E$方向\ $B$方向\ 入口\ 知2求1\quad $V=\dfrac{E}{B}$

只选速度\ 不选电量电性质量（忽略$G$）

若·$V$偏大\ 偏小$\to$螺旋线运动
%%%END
%%%BLOCK id=020-03 ch=11 sec=复合场·质谱仪 kind=method alt=none cont=none y=0.44-0.66
\textbf{质谱仪}\quad $R_{\text{偏}}$越大\ 比荷越小
%FIG p020_d 0.09 0.48 0.47 0.60
\notefig[0.45]{figures/p020_d.png}
\[ qU=\frac12 mv^2 \]
\[ R=\frac{mv}{qB} \]
\[ R=\frac1B\sqrt{\frac{2mU}{q}} \]
\[ R=\frac{mv}{qB}\qquad \frac{q}{m}=\frac{v}{BR}\qquad R\propto\frac{1}{\,q/m\,} \]
%%%END
%%%BLOCK id=020-04 ch=11 sec=复合场·磁流体发电机·电磁流量计·霍尔元件 kind=method alt=none cont=none y=0.66-0.78
\textbf{磁流体发电机\ /\ 电磁流量计\ /\ 霍尔元件}
\[ BqvB=qE=q\frac{U}{d}=q\frac{\varepsilon}{d} \]
% CHECK: 开头手写为 "BqvB"，多出一个 B，照原样转录
\[ \varepsilon=Bdv \qquad v=\frac{U}{Bd} \]
\[ qvB=qE=q\frac{U}{d}\ \leftarrow\ \text{霍尔电压} \]
充电后二力平衡
\[ U=B\,\text{\struck{d}}\,\frac{I}{ne\,\text{\struck{d}}\,h} \]
% 原稿：上式中分子 B 后的 d 与分母中的 d 被斜线划去（约分）
\[ =\frac{1}{neh}\cdot BI\propto\frac1h \]
% 原稿：标题"磁流体发电机/电磁流量计/霍尔元件"中的两条斜线为分隔符；前两行公式在前两项之下，后面的公式在"霍尔元件"之下
%%%END
%%%BLOCK id=020-05 ch=11 sec=复合场·回旋加速器 kind=method alt=none cont=none y=0.78-1.0
\begin{itemize}
\item \textbf{回旋加速器}\quad $T_{\text{交变电场}}=T_{\text{粒子在磁场中}}=\dfrac{2\pi m}{Bq}$
\item 最大速度\quad $V_{\max}=\dfrac{qBR}{m}$
\item 加速次数\quad $\frac12 mV_{\max}^2=nqU$
\item $t_{\text{总}}=t_{\text{电}}+t_{\text{磁}}$
\end{itemize}
% 原稿：以上四行左侧有大括号；"t总=t电+t磁"下一行左端另有一个孤立的"="，其后即下方方框结果
\[ t_{\text{电}}=\frac{qBR}{qE}=\frac{BR}{E} \]
\[ qEt_{\text{电}}=mV_{\max}\ \to\ [\text{当成匀加速直线运动处理}] \]
% 原稿："mV_max"上方标注"动量"
\[ t_{\text{磁}}=\frac{n\pi m}{qB}\qquad (n\to n-1) \]
\begin{keybox}
\[ V_{\max}=\frac{qBR}{m} \]
\[ t=\frac{n\pi m}{qB}+\frac{BR}{E} \]
\end{keybox}
\begin{keybox}
\[ n=\frac{q^2B^2R^2}{q\,2mU}=\frac{qB^2R^2}{2mU} \]
\end{keybox}
%%%END
%%%BLOCK id=021-09 ch=11 sec=洛伦兹力分解·螺旋运动 kind=note alt=12 cont=next y=0.88-0.92
\textbf{螺旋运动}
%%%END
%%%BLOCK id=022-01 ch=11 sec=洛伦兹力分解·动生电动势 kind=derivation alt=12 cont=prev y=0.00-0.125
% 页面顶部被截断，图与上方小字不完整
%FIG p022_a 0.01 0.0 0.205 0.09
\notefig[0.3]{figures/p022_a.png}
（图左下标注）$-BV_yqV_x$

非静电力\unclear{做功}
\[
\text{电动势}\ E=\frac{BV_xqL}{q}=BLV_x
\]
\unclear{总}功　$W_{\text{合}}=0$
%%%END
%%%BLOCK id=036-02 ch=11 sec=复合场·板块起飞临界 kind=method alt=07,06 cont=none y=0.115-0.40
\textbf{2.}\ 运动过程共速，起飞等分开求解

$qEB$ 加杆约束，复合场常出现合力为0、速度不为0的匀速直线运动临界：
\begin{itemize}
\item $qV_0B=F_{\text{外}}$ 时是做匀速运动的临界
\item 或是起飞的临界（即离开木板）
\end{itemize}

%FIG p036_b 0.265 0.158 0.60 0.262
\notefig[0.45]{figures/p036_b.png}

未共速时\unclear{就}起飞
\[
m_2V_2+m_1V_1=(m_1+m_2)V_{\text{共}}
\]

板进了块还\unclear{得算进}时速度\quad $V_{\text{共}}=\cdots$

$qV_0B=m_1g$ 时\quad $V_0>V_{\text{共}}$ 则会共速

\hspace{4.2em}$V_0<V_{\text{共}}$ 则无法共速
\begin{align*}
\therefore\ Q&=\frac12m_1\bigl(V_1^2-V_{1\text{共}}^2\bigr)+\frac12m_2\bigl(V_2^2-\unclear{V_{2\text{共}}^2}\bigr)\\ % 右端被页边切断，第二项括号内后半为推测
&\quad+\mu m_1gS_{\text{相}}
\end{align*}
%%%END
%%%BLOCK id=048-01 ch=11 sec=磁场·提纲 kind=summary alt=none cont=none y=0.09-0.58
\textbf{磁场}
% 页左侧的手写标题"磁场"，其右为下列各条目（各条目独占一行、大段留白）
\begin{itemize}
\item 磁场、磁感应强度
\item 磁感线、电流周围磁场，
\item 安培力\ 及\ 安培力冲量．
\item 安培定则的应用和磁场叠加
\item 安培力的大小和方向
\item 安培力作用下的运动状况
\item 安培力作用下导体的平衡问题．
\item 安培力作用下力电综合问题
\end{itemize}
%%%END
%%%BLOCK id=048-02 ch=11 sec=磁场·提纲 kind=summary alt=none cont=none y=0.68-1.0
\textbf{带电粒子在磁场中的动运．}
% CHECK: 页左侧标题原稿写作"动运"（应为"运动"）；该标题位于"洛伦兹力、洛伦兹力冲量"条目之下、"洛伦兹力的理解与运用"条目之左
\begin{itemize}
\item 洛伦兹力、洛伦兹力冲量
\item 洛伦兹力的理解与运用
\item 粒子在有界磁场中的运动
\item 带电粒子在磁场中的临界极值问题
  \begin{itemize}
  \item 法1、\ 临界条件，边界条件\ 矢量图
  \item 法2．\ 三角函数，二次方程判别式，不等式，图像法
  \end{itemize}
\item 带电粒子在匀强磁场中的多解问题
  \begin{itemize}
  \item 电性不确定、磁场方向不确定、临界状态不\unclear{唯}一（轨迹相切），运动周期性
  \end{itemize}
\end{itemize}
%%%END
%%%BLOCK id=049-01 ch=11 sec=磁场·提纲 kind=summary alt=none cont=none y=0.04-0.98
\red{①}

四类动态圆模型

带电粒子在组合场中的运动

\red{②}

带电粒子在交变电磁场中的运动

\red{③}\ \unclear{E}

带电粒子在叠加场中的运动

① ②
% 页底左侧(y≈0.87 与 0.94)为黑色圈号①②，其后无文字
%%%END
%%%BLOCK id=050-01 ch=11 sec=有界磁场·作图找圆心 kind=method alt=none cont=none y=0.025-0.28
两两垂直、一垂一点作中垂\ \ 直径所对圆周角$90^\circ$

\red{中垂线+垂线}

%FIG p050_b 0.49 0.03 0.87 0.275
\notefig[0.55]{figures/p050_b.png}
% 红笔所绘：一个大红圆，圆内有竖直双线(中垂线)、水平弦及斜线构成的三角形，右侧有红色箭头；叠在"磁场中平抛"一节右上部的公式之上，公式见块 050-02
%%%END
%%%BLOCK id=050-02 ch=11 sec=组合场·磁场中平抛 kind=derivation alt=04,09 cont=none y=0.06-0.31
\fcolorbox{red!75!black}{white}{\textbf{磁场中平抛}．}

%FIG p050_a 0.03 0.100 0.315 0.25
\notefig[0.4]{figures/p050_a.png}
% 图内标注：V_0、E(两条竖直向下箭头)、y、x、2α(红)、θ(红圈)、V_1、B(虚线圆内为点，表示磁场)

\[
\begin{cases}\tan\theta=\dfrac{2y}{x}=2\tan\alpha=\dfrac{V_y}{V_0}\\[6pt] V_1\cos\theta=V_0\end{cases}
\qquad
\begin{cases}x=V_0t\\[2pt] y=\dfrac{1}{2}\dfrac{qE}{m}t_{\theta}^{2}\\[4pt] t_{\theta}=\dfrac{V_y}{a}\end{cases}
\]
% CHECK: 右侧大括号中 t 的下标被红笔覆盖，字形像 θ 或 e，按 t_θ 录入

洛伦兹力提供向心力

\[
\begin{cases}qV_1B=\dfrac{mV_1^{2}}{R}\\[6pt] T=\dfrac{2\pi R}{V_1}\\[6pt] \dfrac{t}{T}=\dfrac{\theta}{2\pi}\ \ \text{or}\ \ t=\dfrac{\theta R}{V_1}\ \ \text{（无}\,B\,\text{求时间）}\end{cases}
\qquad
\begin{cases}R=\dfrac{mV}{qB}\\[6pt] T=\dfrac{2\pi m}{qB}\\[6pt] t=\dfrac{\theta m}{qB}\end{cases}
\]

\[ D=\frac{2mV}{qB}\sin\theta=\frac{2mV_y}{qB}=\frac{m\Delta V_y}{qB} \]

\[ a_{\text{向}}=\omega V \]

☆\ \fcolorbox{red!75!black}{white}{$qBD=m\Delta V_y$}
\[ \Delta V_y=2V_y=2at=\frac{2Eqx}{mV_0} \]
\[ D=\frac{2Ex}{BV_0}\ \ \text{（与粒子无关）} \]
% "（与粒子无关）"原稿为涂改：先写"看"，上方改写"与"，并被圈住

（正负\unclear{加}力量，与侧移结合）

侧移量只与场\unclear{强}有关
%%%END
%%%BLOCK id=050-03 ch=11 sec=有界磁场·作图找圆心 kind=method alt=none cont=none y=0.285-0.43
速偏$=$圆心$=2\cdot$位偏

几何长$\to$\fbox{半径}$\to$其他量

%FIG p050_c 0.01 0.333 0.245 0.42
\notefig[0.4]{figures/p050_c.png}
% 图内含(黑字)"有垂直即相等""(垂径定理)"及红笔标记

有垂直即\red{相等}\quad （垂径定理）
% 原稿"有垂直即"之后黑字被红笔"相等"叠写覆盖

%FIG p050_d 0.24 0.332 0.415 0.425
\notefig[0.3]{figures/p050_d.png}
% 图内标注：V、D、L、r、r−D、×(磁场区)；红笔另画有 L 形大括线及 D、L 标记

双边限制\quad \fcolorbox{red!75!black}{white}{$\displaystyle r=\frac{D^{2}+L^{2}}{2D}=\frac{mV}{qB}$}

%FIG p050_e 0.585 0.335 0.77 0.397
\notefig[0.3]{figures/p050_e.png}
2种临界
%%%END
%%%BLOCK id=050-04 ch=11 sec=组合场·电场与磁场偏转对比 kind=method alt=04,09 cont=none y=0.40-0.51
%FIG p050_f 0.01 0.415 0.115 0.495
\notefig[0.25]{figures/p050_f.png}
% 图内标注：V_0、y、x

\[
\left\{\begin{array}{l}
\text{平抛}\ \ \fcolorbox{red!75!black}{white}{$\displaystyle y=\frac{g}{2V_0^{2}}x^{2}$}\\[10pt]
\fcolorbox{red!75!black}{white}{圆周}\ \ r=\dfrac{mV_0}{qB}=\fcolorbox{red!75!black}{white}{$\displaystyle\frac{x^{2}+y^{2}}{2y}$}
\end{array}\right.
\]
% "平抛"与"y="之间另有一个红色小记号(像 ∝ 或 ω)

平进斜出

%FIG p050_g 0.335 0.42 0.53 0.50
\notefig[0.35]{figures/p050_g.png}
% 图内标注：L、d/2、d/2、纯E、纯B（红笔在图上画有轨迹线与圈）

\[ \frac{d}{2}=\frac{1}{2}\,\frac{qE}{mV_0^{2}}L^{2}\qquad qE=\frac{dmV_0^{2}}{L^{2}} \]
\[ \frac{mV}{qB}=\frac{\left(\frac{d}{2}\right)^{2}+L^{2}}{2\cdot\frac{d}{2}}\qquad qV_0B=\frac{dmV_0^{2}}{L^{2}+\frac{d^{2}}{4}} \]
\fcolorbox{red!75!black}{white}{\begin{tabular}{l}故 $qE>qV_0B$\\ 电场力大于洛伦兹力\end{tabular}}
%%%END
%%%BLOCK id=050-05 ch=11 sec=有界磁场·圆形边界 kind=method alt=none cont=none y=0.505-0.71
指心入背心出\quad 双心连线．三点共线\quad \fcolorbox{red!75!black}{white}{速交双心}
% 红笔在"指心入背心出""双心连线""三点共线"下各画有一道红线

\unclear{夹手}指角入\ 背\unclear{心}角出
% CHECK: 此行字迹潦草，参照上一行"指心入背心出"，疑为"夹角指心角入、背心角出"一类，待核；"背""角"之间下方另有一个"夫/关"形小字

%FIG p050_h 0.01 0.543 0.19 0.685
\notefig[0.3]{figures/p050_h.png}
% 图内标注：夹角(?)、L、r、θ，红笔画有若干弧线与标记

\[
\begin{cases}D=(L+R)\tan\theta\\[4pt] \tan\dfrac{\theta}{2}=\dfrac{R}{r}\\[8pt] \tan\theta=\dfrac{2\tan\frac{\theta}{2}}{1-\tan^{2}\frac{\theta}{2}}\end{cases}
\]

%FIG p050_i 0.39 0.566 0.635 0.715
\notefig[0.4]{figures/p050_i.png}
% 图内标注：R、r、M 等；红笔画有若干大圈与竖线
%%%END
%%%BLOCK id=050-06 ch=11 sec=有界磁场·直线边界 kind=method alt=none cont=none y=0.545-0.75
%FIG p050_j 0.595 0.545 0.975 0.677
\notefig[0.6]{figures/p050_j.png}
% 图内标注：正电、负电、Z、D、φ（多处）、V_y=V sinφ、π−φ、"侧移相等"；红笔圈出 V_y=V sinφ

侧移相等

\[ D=2r\sin\varphi=\frac{2mV\sin\varphi}{qB} \]
\fcolorbox{red!75!black}{white}{正负\unclear{加}力量}
\[ t_1+t_2=T=\frac{2\pi m}{qB} \]
% CHECK: "正负加力量"中第三字字形不清(像"加"或"后"，上方两处同样出现此四字)，含义待核
%%%END
%%%BLOCK id=050-07 ch=11 sec=有界磁场·相切临界 kind=method alt=none cont=none y=0.727-0.98
%FIG p050_k 0.01 0.727 0.255 0.937
\notefig[0.4]{figures/p050_k.png}
% 图内含标题"直线相切"，及标注：切点、角平分线、虚线(延长线)等

1延\ \ 2交\ \ 3分\ \ 4垂\ \ 5再垂
% 红笔在"1延""2交""4垂"下各画一条红线
%%%END
%%%BLOCK id=050-08 ch=11 sec=有界磁场·相切临界 kind=method alt=none cont=none y=0.727-0.98
%FIG p050_l 0.285 0.727 0.535 0.937
\notefig[0.4]{figures/p050_l.png}
% 图内含标题"圆形相切""束缚粒子"，及标注：r、R、2R−r 等（同心大小圆）

\[ \fcolorbox{red!75!black}{white}{$\displaystyle (2R-r)^{2}=r^{2}+R^{2}+2rR\,\overset{\text{补角}}{\underline{\cos\theta}}$} \]

1垂\ \ 2射\ \ 3等\ \ 4交\ \ 5再交
%%%END
%%%BLOCK id=050-09 ch=11 sec=有界磁场·极值 kind=method alt=none cont=none y=0.86-0.95
直径所对圆周角为直角

$V_{\min}$\quad 位移/\unclear{距离}当直径

$V_{\max}$\quad 最先射出找$R_{\max}$
% CHECK: "位移"与"距离"之间的"/"原稿为一条斜线，含义(或为"位移(距离)当直径")待核
%%%END
%%%BLOCK id=051-01 ch=11 sec=复合场·磁流体发电机 kind=knowledge alt=10 cont=none y=0.03-0.30
\textbf{磁流体发电机}（原稿圈出标题）

%FIG p051_a 0.10 0.10 0.33 0.19
\notefig[0.3]{figures/p051_a.png}

\[
\left\{\begin{aligned}
&qvB=qE\\
&U=Ed=Bdv
\end{aligned}\right.
\]
\[
I=\frac{U}{R+r}=\frac{Bdv}{R+\dfrac{\rho d}{S_{\text{板}}}}
\]
\[
\Delta P_{\text{压}}S=F_{\text{安}}\,\unclear{+f}=BId+f
\]
\[
\eta=\frac{Bdv\cdot I}{\Delta P_{\text{压}}S\cdot v}=\frac{UI}{Fv}=\frac{UI}{(F_{\text{安}}+f)v}
\]
% 上一行第二项分母原稿为 F v（未写下标）

$P_{\text{机}}\to P_{\text{电}}\to$ \unclear{阻} $\to P_{\text{\unclear{内}}}$
%%%END
%%%BLOCK id=051-02 ch=11 sec=复合场·电磁流量计 kind=knowledge alt=none cont=none y=0.26-0.42
\textbf{电磁流量计}（原稿圈出标题）

%FIG p051_b 0.09 0.33 0.243 0.385
\notefig[0.25]{figures/p051_b.png}

\[
Q=\frac{V_{\text{体积}}}{t}=S\cdot v=ac\,\frac{U_0}{Ba}
\]

%FIG p051_c 0.50 0.33 0.64 0.42
\notefig[0.25]{figures/p051_c.png}
%%%END
%%%BLOCK id=051-03 ch=11 sec=复合场·霍尔效应 kind=knowledge alt=none cont=none y=0.38-0.55
\textbf{霍尔效应}（原稿圈出标题）

%FIG p051_d 0.06 0.435 0.43 0.53
\notefig[0.4]{figures/p051_d.png}

\[
\left\{\begin{aligned}
&U_H=Bdv=\frac{1}{ne}\cdot\frac{BdI}{S_{\text{\unclear{…}}}}=R_H\,\frac{BdI}{S}\qquad \text{$R_H$为霍尔系数}\\
&I=nesv
\end{aligned}\right.
\]
% 第一个分式分母 S 右下角有不明小标记
%%%END
%%%BLOCK id=051-04 ch=11 sec=复合场·回旋加速器 kind=knowledge alt=09 cont=none y=0.52-0.97
\textbf{回旋加速器}（原稿圈出标题）

%FIG p051_e 0.095 0.592 0.32 0.713
\notefig[0.25]{figures/p051_e.png}

\begin{align*}
v_m&=\frac{qBR}{m} & E_{km}&=\frac{q^2B^2R^2}{2m} & P_m&=qBR\\
n&=\frac{E_{km}}{qU}=\frac{qB^2R^2}{2mU}\\
t_B&=\frac{n\pi m}{qB}=\frac{\pi R\,\unclear{n}}{v_m}\\
t_E&=\sqrt{\frac{2nmd^2}{qU}}=\frac{2nd}{v_m}\qquad nd=\frac12\,\frac{qU}{md}\,t_E^2
\end{align*}
由 $t_B\gg t_E$：$\dfrac{\pi}{2}$、$\dfrac{R}{d}\ge1$

\[
v=2\pi Rf\qquad E_k=\frac12mv^2
\]
\[
\frac{qBR}{m}=2\pi Rf
\]
\[
f=\frac{qB}{2\pi m}\qquad \text{$m$不变}\quad B\times2\quad f\,\text{\unclear{(↑)}}\times2
\]

$f=\dfrac{qB}{2\pi m}$（交变电压频率），\quad $f\propto B$ \quad B最大为$B_m$ \quad $f_m=\dfrac{qB_m}{2\pi m}$
\[
\left\{\begin{aligned}
&\underset{\text{\scriptsize 实际}}{f_m}<\underset{\text{\scriptsize 标准}}{f_m}\ \text{时}\quad v_m=2\pi f_mR\\
&\underset{\text{\scriptsize 实际}}{f_m}>\underset{\text{\scriptsize 标准}}{f_m}\ \text{时}\quad v_m=\frac{qBR}{m}
\end{aligned}\right.
\]
% 第一式 v_m= 之后、2π 之前有涂改痕迹；第二式分子 qBR 有涂改
% 手写 v 与 V 不分，这里一律按速度记为小写 v（v_m 原稿写作 V_m）
%%%END
%%%BLOCK id=083-03 ch=11 sec=磁感应强度·定义式 kind=formula alt=none cont=none y=0.07-0.12
\[
B=\frac{F}{IL\sin\theta}
\]
%%%END
%%%BLOCK id=083-10 ch=11 sec=磁场·受力条件与方向 kind=note alt=none cont=none y=0.24-0.37
% 原稿此块及下块左侧有红笔大括号
在$B$中放$I$\quad 测不了$B$\quad $B=\dfrac{F}{IL\sin\theta}$\quad $\theta$未知.

$I$不受磁场力\quad ①$B=0$\quad ②$I\parallel B$\qquad 粒子：①$B=0$\quad ②$v\parallel B$

电流受磁场力方向与磁场垂直\qquad 三垂直\quad $F=qv\times B$\quad $F=IL\times B$

\unclear{场}真\quad 线假
%%%END
%%%BLOCK id=089-04 ch=11 sec=带电粒子在磁场中的圆周·动量 kind=problem alt=07 cont=none y=0.37-0.51
\noindent 动量+磁场
\[ qVB=\frac{mV^2}{R}\qquad R=\frac{mV}{qB}=\frac{P}{qB} \]
%FIG p089_c 0.41 0.385 0.53 0.445
\notefig[0.25]{figures/p089_c.png}
\[ m_1V_0=(m_1+m_2)V_{\text{共}} \]
\[ r_1=\frac{m_1V_0}{qB}\qquad r_{\text{共}}=\frac{(m_1+m_2)V_{\text{共}}}{qB}=\frac{m_1V_0}{qB}=r_1 \]
\noindent 由 $V_{\text{共}}<V_0$ 故用时更长。
%%%END
%%%BLOCK id=090-01 ch=11 sec=有界磁场·几何关系 kind=method alt=none cont=none y=0.03-0.185
\noindent \textbf{几何关系处理}
% 页首标题；本块与下一块为圆弧轨迹的几何关系（原稿未写 B、q，按图形推断为有界磁场中的圆周）

%FIG p090_a 0.09 0.055 0.385 0.175
\notefig[0.5]{figures/p090_a.png}
\[ \begin{cases} R\cos\alpha+R\sin\alpha=a\\ R-R\sin\alpha=\dfrac a2\\ \sin^2\alpha+\cos^2\alpha=1\end{cases}\qquad R=\left(\dfrac{4-\sqrt6}{2}\right)a \]
% 手写中 a 的字形近似 9
%%%END
%%%BLOCK id=090-02 ch=11 sec=有界磁场·几何关系 kind=method alt=none cont=none y=0.185-0.31
%FIG p090_b 0.06 0.18 0.27 0.305
\notefig[0.3]{figures/p090_b.png}
\[ r^2+R^2-2rR\cos\varphi=\left[r-(d-R)\right]^2 \]
% 括号内手写形似 “r(d−R)”，按图中底边标注取 r−(d−R)
%%%END
%%%BLOCK id=091-03 ch=11 sec=有界磁场·极值 kind=method alt=90 cont=none y=0.29-0.34
2020选择7中\unclear{磁}场 $\theta$ 也可\underline{求导求最小}：
\[
\cos\theta=\frac{r^{2}+4R^{2}-4Rr+r^{2}-R^{2}}{2(2R-r)r}\qquad
\cos'\theta=\frac{12(r-R)R^{2}}{[2(2R-r)r]^{2}}
\]
% CHECK: 分母括号内 2R 与 r 之间的减号手写模糊，按上式取 2R-r
%%%END
%%%BLOCK id=092-03 ch=11 sec=质谱仪·狭缝 kind=method alt=none cont=none y=0.42-0.61
①\ \textbf{质谱仪狭缝}
%FIG p092_c 0.055 0.515 0.215 0.605
\notefig[0.25]{figures/p092_c.png}
\[
\left\{\begin{aligned}
d&\le 2r_{1}-2r_{2}\\
L-d&=2r_{2}
\end{aligned}\right.
\qquad d\le\frac{\sqrt{m_{1}}-\sqrt{m_{2}}}{2\sqrt{m_{1}}-\sqrt{m_{2}}}\,L
\qquad \frac{r_{1}}{r_{2}}=\sqrt{\frac{m_{1}}{m_{2}}}
\]
% CHECK: d≤ 结果式的分母手写为 2√m1−√m2，请核对
%%%END
%%%BLOCK id=092-04 ch=11 sec=质谱仪·狭缝 kind=method alt=none cont=none y=0.61-0.90
②\ 轨迹不交叠 $\Leftrightarrow r_{1}\ge r_{2}$
% CHECK: “不交叠”后的符号手写像 ⇔（也可能是 ⇒）
%FIG p092_d 0.13 0.658 0.40 0.805
\notefig[0.3]{figures/p092_d.png}
\[
\left\{\begin{aligned}
d&=r_{1}-h\\
h&=\sqrt{r_{1}^{2}-\frac{L^{2}}{4}}
\end{aligned}\right.
\]
\[
\therefore\ d=\frac{2}{B}\sqrt{\frac{mU_{0}}{q}}-\sqrt{\frac{4mU_{0}}{qB^{2}}-\frac{L^{2}}{4}}
\]
%%%END
%%%BLOCK id=092-05 ch=11 sec=质谱仪·底片重叠条件 kind=problem alt=none cont=none y=0.73-0.99
若两离子\unclear{在}底片上没有重叠\ 求 $L$ 条件
\begin{align*}
r_{\text{甲}}&=\sqrt{\frac{4m(U-\Delta U)}{qB^{2}}}\\
r_{\text{乙}}&=\sqrt{\frac{2m(U+\Delta U)}{qB^{2}}}\\
2r_{\text{甲}}-2r_{\text{乙}}&\ge L
\end{align*}
% CHECK: r乙 的下标手写像 2（乙），ΔU 的 Δ 手写像 δ
甲最近点$-$乙最远点$\ge L$
\[
L\le2\sqrt{\frac{4m(U-\Delta U)}{qB^{2}}}-2\sqrt{\frac{2m(U+\Delta U)}{qB^{2}}}
\]
%%%END
%%%BLOCK id=097-01 ch=11 sec=磁感应强度·安培力 kind=formula alt=none cont=none y=0.03-0.17
磁场\quad $B=\dfrac{F}{IL\sin\theta}=\dfrac{F}{qv\sin\theta}$\qquad $B=\dfrac{kI}{r}$\quad 通电导线间\unclear{安培}力；导线\unclear{同}吸异斥。（不考虑有外磁场）

%FIG p097_a 0.13 0.10 0.22 0.17
\notefig[0.25]{figures/p097_a.png}
1连\quad 2垂\quad 3右手

联系：$F_{\text{安}}=BIL=B\,nesv\,L=NBev=NqvB$
%%%END
%%%BLOCK id=097-02 ch=11 sec=通电导线间作用力·技巧 kind=method alt=none cont=none y=0.03-0.12
技巧\ $\otimes$\qquad 直接画力的\unclear{…}\\
跳过画$B$
% 原稿右侧被页边截断，“直接画力的”之后的字不可见
%FIG p097_b 0.77 0.03 1.0 0.12
\notefig[0.25]{figures/p097_b.png}
%%%END
%%%BLOCK id=097-03 ch=11 sec=安培力·有效长度 kind=knowledge alt=none cont=none y=0.19-0.22
有效长度连首尾，$\therefore$有效长度越大越易失去平衡。
%%%END
%%%BLOCK id=098-01 ch=11 sec=洛伦兹力分解·速度分量 kind=diagram alt=06 cont=none y=0.07-0.23
①
%FIG p098_a 0.03 0.08 0.21 0.20
\notefig[0.25]{figures/p098_a.png}

$E_{k}\uparrow$\quad 支持力做功

合外力为0时\ $V_{\max}$
% 原稿“合外”二字被圈起
%%%END
%%%BLOCK id=098-05 ch=11 sec=洛伦兹力·动量定理 kind=problem alt=09,03,07 cont=none y=0.58-0.77
④
%FIG p098_d 0.02 0.587 0.225 0.695
\notefig[0.25]{figures/p098_d.png}

\unclear{$t_s$}离开地面，求位移
\[ \begin{cases} Eqt-\mu(mg-qvB)t=mV\\[2pt] qvB=mg \end{cases} \qquad x=\frac{\mu mgt}{qB} \]
% CHECK: x 的分子分母手写如此（μ 似应在分母中同样出现而相消）

$Eqt=mV$
\[ a=\frac{Eq}{m}-\mu g+\frac{qVB}{m} \]
% CHECK: 末项手写未见 μ（按受力似应为 μqVB/m）
\[ V=\frac{Eqt}{m}-\mu gt+\frac{qBx}{m} \]
% CHECK: 末项手写未见 μ（似应为 μqBx/m）

一个动量定理\quad $V=\dfrac{mg}{qB}$
%FIG p098_e 0.672 0.585 0.815 0.658
\notefig[0.25]{figures/p098_e.png}
%%%END
%%%BLOCK id=114-02 ch=11 sec=洛伦兹力·圆周运动 kind=problem alt=04 cont=none y=0.25-0.46
\begin{problem}[2]
% 图p114_b：N、qvB、mg 的力三角形（含θ）；图p114_c：球在碗内表面运动的受力图（N、qvB、mg）
%FIG p114_b 0.225 0.248 0.35 0.31
\notefig[0.25]{figures/p114_b.png}

%FIG p114_c 0.01 0.305 0.19 0.40
\notefig[0.3]{figures/p114_c.png}

球能圆周运动 求 $B_{\min}$ 及对应速率
\begin{solution}
【法1】\quad $\dfrac{mv^2}{r}-qvB+mg\tan\theta=0$\qquad $\Delta\ge0\ \Rightarrow\ B$

【法2】\quad $\left\{\begin{array}{l}qvB-N\sin\theta=\dfrac{mv^2}{r}\\[2mm] N\cos\theta=mg\end{array}\right.$

$\therefore\ B=\dfrac{mv}{qr}+\dfrac{mg\tan\theta}{qv}\ge2\sqrt{\dfrac{m^2g\tan\theta}{q^2r}}$ % 此行右端被拍摄边缘截断（可见 m^2 g tan…/q^2 r），按可见部分转录

当 $\dfrac{mv}{qr}=\dfrac{mg\tan\theta}{qv}$ 时取等
\end{solution}
\end{problem}
%%%END
%%%BLOCK id=127-01 ch=11 sec=洛伦兹力·杆上带电小环 kind=method alt=03 cont=none y=0.02-0.30
% 图内手写（红字，辨认不全）：f_摩（沿杆向上的箭头）、F_洛（垂直杆斜向上的箭头）、N（垂直杆斜向下的箭头，旁有“先/后”字样）、mg（竖直向下）；小环带“+”，×表示磁场垂直纸面向里
%FIG p127_a 0.005 0.02 0.31 0.168
\notefig[0.3]{figures/p127_a.png}
\[ a_1\,\red{\uparrow}=\dfrac{mg\sin\theta-\unclear{\mu}\,f\downarrow}{m} \]
% CHECK: 减号后的“μ”像被划去/写重了，且 f 已含 μ；原样转录
% 图内红字：f=0时、N=0（二者之间有下划线）；a_max（红字，标在曲线拐点附近）
%FIG p127_b 0.37 0.029 0.67 0.24
\notefig[0.3]{figures/p127_b.png}
\[ a_2=\dfrac{mg\sin\theta-f\uparrow}{m}\ \downarrow\qquad \rule[0.5ex]{3em}{0.4pt}\ v_{\max} \]
\[ \left\{\begin{array}{l} a_{\max}\text{时}\ f,N=0\\[2pt] v_{\max}\text{时}\ \Sigma F=0 \end{array}\right. \]
%%%END
%%%BLOCK id=131-06 ch=11 sec=地磁场 kind=note alt=none cont=none y=0.62-0.86
% 页面右侧，周围有大片波浪线涂划；“地磁南北是”下方有一道波浪线；其上方另有几个像“m m”的涂写符号
地磁南北是\ 反向、

不是匀强磁场
%%%END
%%%BLOCK id=150-06 ch=11 sec=磁感应强度·定义式 kind=formula alt=none cont=none y=0.72-0.75
\noindent $B=\dfrac{F}{IL\sin\theta}$　　电生磁　一\unclear{速}二\unclear{垂}三右手% “速”“垂”两字笔迹难辨
%%%END
%%%BLOCK id=150-07 ch=11 sec=洛伦兹力分解·分力做功 kind=problem alt=06 cont=none y=0.75-0.90
% 图中红笔标注：F洛1、V1、V合、F洛2 等；黑笔有 qVB、V、mg；图为地面上的小木块，其周围画有“×”(磁场垂直纸面向里)
\noindent 洛伦兹力　分力可以做功，合力不做功

\begin{problem}
%FIG p150_i 0.045 0.776 0.225 0.88
\notefig[0.25]{figures/p150_i.png}
使小木块对地分离　求磁场移动速度
\begin{solution}
$qVB=mg$　使$B$向左以$V=\dfrac{mg}{qB}$运动
\end{solution}
\end{problem}
%%%END
%%%BLOCK id=158-04 ch=11 sec=复合场·磁流体发电机 kind=problem alt=10,06 cont=none y=0.60-1.00
% 图：矩形管道，标有 R(与开关串联)、B、V0(流速箭头)、d、L、h，图内文字已在图中
\begin{problem}
%FIG p158_f 0.04 0.585 0.262 0.752
\notefig[0.25]{figures/p158_f.png}
\noindent 液流，电阻率$\rho$，匀速$V_0$

\noindent ①　$qV_0B=\dfrac{qU_0}{d}$　$\therefore U_0=BdV_0$

\noindent ②　开关闭合前后管道两端压强差.
\[ \left\{\begin{aligned}
P_1hd&=f\\
P_2hd&=f+f_{\text{安}}\\
F_{\text{安}}&=BIL\quad I=\dfrac{U_0}{R+\dfrac{\rho d}{Lh}}
\end{aligned}\right.\qquad \Delta P=\dfrac{LdV_0B^2}{LhR+d\rho} \]% CHECK: F_安=BIL 末字原稿形似L，按 ΔP 结果(含d)应为 d

\noindent ③　调整宽和高，使$S=dh$不变，求$P_{R\max}$及此时$\dfrac{d}{h}$值
\[ P=\dfrac{(BLV_0)^2}{R+\dfrac{\rho d}{Lh}}\qquad P_{\max}=\dfrac{Lh(BLV_0)^2}{2\rho d} \]% CHECK: (BLV_0) 中的 L 原稿字形似 L，按 U_0=BdV_0 应为 d；P_max 分母原稿为 2ρd(按 U_0^2/4r 推算应为 4ρd)

\noindent 当$R=\dfrac{\rho d}{Lh}$时$P_{\max}$　$\dfrac{d}{h}=\dfrac{LR}{\rho}$
\end{problem}
%%%END
%%%BLOCK id=161-01 ch=11 sec=摆线运动·最大速度 kind=derivation alt=none cont=none y=0.04-0.25
$V_x=ky$ 与 $E$ 无关，求 $V_{\max}$

%FIG p161_a 0.31 0.04 0.60 0.145
\notefig[0.3]{figures/p161_a.png}
\[ qEy_m=\tfrac12mV_m^2-\tfrac12mV_0^2 \]
$V_m=ky_m$\quad \unclear{令}$E=0$\quad 又$V_0=kR_0$\quad 故 \struck{$V_m=kR_0$}\quad $y_0=R_0$

% 下一行 k=(q/m)B 之后有涂改液痕迹
\[ k=\frac{q}{m}B\qquad V_m^2-\frac{2E}{B}V_m-V_0^2=0\qquad V_m=\frac{E}{B}\pm\sqrt{\left(\frac{E}{B}\right)^2+V_0^2} \]
%%%END
%%%BLOCK id=161-02 ch=11 sec=摆线运动·配速法 kind=method alt=none cont=none y=0.25-0.59
\textbf{〈配速法〉}

%FIG p161_b 0.13 0.27 0.40 0.405
\notefig[0.3]{figures/p161_b.png}
\[ V_0=\frac{mg}{qB}\qquad 0=1+(-1). \]
\[ \text{两个分运动}\left\{\begin{array}{l}\text{向右匀速直线运动}\\ \text{逆时针圆周运动}\end{array}\right.\ \to\ \text{摆线} \]

\textbf{轨迹为摆线}
\[ V_{\max}=V_0+V_0=\frac{2mg}{qB}\qquad V_{\min}=V_0-V_0=0 \]
竖直最大位移\quad$mgh=\tfrac12mV_{\max}^2$\quad$h=\dfrac{2m^2g}{q^2B^2}$\quad$R=\dfrac{mV_0}{qB}=\dfrac{m^2g}{q^2B^2}$\quad$\therefore h=2R$ 为直径长度

$T=\dfrac{2\pi m}{qB}$\quad 到最低点时间\quad$t=\dfrac{T}{2}=\dfrac{\pi m}{qB}$

1个周期内前进距离\quad$S=V_0T=\dfrac{2\pi m^2g}{q^2B^2}$
%%%END
%%%BLOCK id=161-03 ch=11 sec=摆线运动·洛伦兹力分量冲量 kind=derivation alt=none cont=none y=0.575-0.70
%FIG p161_c 0.10 0.575 0.285 0.655
\notefig[0.25]{figures/p161_c.png}
\[
mgt\ \left\{\begin{aligned}
&mgh=\tfrac12mV_2^2-\tfrac12mV_1^2\\
&qB\Delta x=m\Delta v_y\\
&qB\Delta y\,\text{\struck{$+mgt$}}=m\Delta v_x
\end{aligned}\right.
\]
〈洛伦兹分力冲\unclear{量}〉 % 末字上有涂改
%%%END
%%%BLOCK id=161-04 ch=11 sec=有界磁场·圆约束 kind=method alt=none cont=none y=0.71-0.80
% 原稿此块被大圈（云形曲线）圈起，左上角“圆约束”三字另有方框
\begin{keybox}[圆约束]
两次进磁区圆心\unclear{角}$150^\circ$，经$m$次偏转后再入射
\[ m\cdot150^\circ=n\cdot360^\circ \]
\[ m=12\qquad n=5\qquad \text{求最小正整数} \]
\end{keybox}
%%%END
%%%BLOCK id=161-05 ch=11 sec=摆线运动·综合题 kind=problem alt=06,09 cont=none y=0.77-1.00
\begin{problem}
%FIG p161_d 0.01 0.76 0.30 0.955
\notefig[0.3]{figures/p161_d.png}
静止释放，在最低点曲率半径为$2y_m$ % 此行“静止释放”处有涂改液痕迹

① 求小球运动$P(x,y)$速度
\begin{solution}
\[ mgy=\tfrac12mv^2\qquad v=\sqrt{2gh} \]
\end{solution}
② 求第一次下降的最大距离$y_m$
\begin{solution}
\[ qvB-mg=\frac{mv^2}{R}\qquad v_1=\sqrt{2gy_m} \]
\[ R=2y_m\qquad \therefore y_m=\frac{2m^2g}{q^2B^2} \]
\end{solution}
③ 在竖直方向加一匀强电场$E$（$E>\dfrac{mg}{q}$），求$V_{\max}$ % 原稿“匀强”二字插写在“一”字上方
\begin{solution}
\[ (qE-mg)y_m=\tfrac12mV_m^2 \]
\[ qV_mB\ \unclear{\pm}\ mg+qE=\frac{mV_m^2}{R}\qquad R=2y_m \] % CHECK: 此式中间的符号有涂改（像 + 与 - 叠写）；按解得结果应为 qV_mB+mg-qE=mV_m^2/R
解得 $V_m=\dfrac{2}{qB}(qE-mg)$
\end{solution}
\end{problem}
%%%END
%%%BLOCK id=163-06 ch=11 sec=安培力·圆环张力 kind=derivation alt=12 cont=none y=0.77-0.87
% 图上标注“超导线圈”、F、BI2R
%FIG p163_f 0.04 0.77 0.265 0.87
\notefig[0.3]{figures/p163_f.png}
\[ 2F_{\text{\unclear{张}}}=BI\,2R\qquad F=BIR \]
%%%END
%%%BLOCK id=167-01 ch=11 sec=带电粒子在复合场中的运动 kind=summary alt=09,04 cont=none y=0.08-0.55
\textbf{粒子运动}

① 电场$+$重力场\quad 等效重力场
\begin{itemize}
\item 类平抛/斜抛
\item 圆周
\item 直线 加减速
\end{itemize}

② 纯磁场
\begin{itemize}
\item $B\perp v$\quad 匀速圆周
\item $B\parallel v$\quad 匀速直线运动
\item $B$ 与 $v$ 夹角 $\ne 90^\circ$，$\ne 0^\circ$\quad 既有垂直分量又有水平分量\quad 等距螺旋线运动
\end{itemize}
%FIG p167_a 0.52 0.198 0.785 0.265
\notefig[0.4]{figures/p167_a.png}
螺旋半径
\[
R=\frac{mv_{\perp}}{qB}=\frac{mv\sin\theta}{qB}
\]
\[
T=\frac{2\pi m}{qB}\quad\text{与 $v$ 无关}
\]
\[
S=\frac{2\pi mv\cos\theta}{qB}=v_{\parallel}T
\]

③ 电磁组合拼接场\quad 速度关联、衔接，形状几何关系

④ 电磁复合叠加场\quad 分析\quad 无约束：直线、圆周、摆线。有约束：$N$ 方向会改变 $\to$ 具体分析

（右上方小注）直接分类讨论 $F$ 和 $qv_0B$ 的关系
% 上行为④行右端上方的小字注记；“的关系”四字写得更小，位于 qv_0B 下方
\par\smallskip
$\hookrightarrow$（自“无约束”一行引出）
\begin{itemize}
\item 直线必匀直
\item 圆周必匀圆
\item 周期由场定，与比荷无关
\end{itemize}
$T=\dfrac{2\pi m}{qB}=\dfrac{2\pi E}{gB}$\qquad $mg=qE$\qquad $F_{\text{洛}}=qvB=\dfrac{mv^2}{R}$
% 上述公式写在括号右侧，并有一斜线箭头指向“周期由场定，与比荷无关”

确定 $B$、$E$ 后，电性、运动方向也确定

⑤ 交变场
\begin{itemize}
\item $E$ 交变\quad $v$-$t$ 图
\item $B$ 交变\quad 按 $B$-$t$ 图画轨迹
\end{itemize}
由于 $B$、$E$ 有周期性，故粒子运动轨迹也有周期性、对称性、重复性
%FIG p167_b 0.355 0.466 0.52 0.545
\notefig[0.3]{figures/p167_b.png}
\[
\tan\theta=\frac{Eq}{mg}
\]
\[
\cos\theta=\frac{mg}{qvB}
\]
%%%END
%%%BLOCK id=167-02 ch=11 sec=带电粒子在复合场中的运动 kind=note alt=06,07 cont=none y=0.55-0.62
电磁场\ 能量\ 动量、力电综合，既有典型的力学过程，又有粒子运动的特点\\
\hspace*{2em}受力、运动、能量、动量
%%%END
%%%BLOCK id=167-03 ch=11 sec=有界磁场·找圆心 kind=method alt=none cont=none y=0.63-0.74
\textbf{找圆心}
\begin{itemize}
\item 出入速度作垂线 % 原稿此处先写“做”，其上方另写了“作”字（疑为改字）
\item 弦作中垂线 % “作”字处有涂改贴/白底
\item 偏转角作角平分线
\item 相切边界作垂线
\end{itemize}
%%%END
%%%BLOCK id=167-04 ch=11 sec=有界磁场·转筒类问题 kind=method alt=none cont=none y=0.76-0.81
转筒类问题\quad $\omega_{\text{粒子}}=\omega_{\text{筒}}$ 为不碰撞条件 % CHECK: 原稿“ω”字形近似“W”，据“转筒”语境读作角速度 ω
%%%END
%%%BLOCK id=169-07 ch=11 sec=安培力·磁电式电流表 kind=knowledge alt=10 cont=none y=0.87-0.99
磁电式电流表（$B$ 不是匀强 而均匀分布）

\underline{\unclear{转}向 同 $F_{\text{安}}$ 方向}，$B$ 不变

铝优点：不会被磁化\quad 减小对磁场影响

$\rho$ 小，电磁阻尼效果好 % “磁”字为后补，小字写在“电”“阻”之间的上方

运输时将接线柱用导线相连防损坏 % 写在页面右下；“运输时将接线柱用导线”下有划线，“相连防损坏”被框起
%%%END
%%%BLOCK id=208-01 ch=11 sec=磁场·有界磁场中粒子偏转与磁场边界 kind=problem alt=04 cont=none y=0.03-0.50
% 题干为粘贴的印刷纸条；页面左缘有一个对勾标记；纸条背面透出的反字已忽略
\begin{problem}[25.（19分）]
如图，在$xOy$平面坐标系的第一象限内有一边长为$a$的正方形区域$OPMN$，该区域内存在垂直$xOy$平面向外的匀强磁场，在$O$点存在一粒子源，可向平面内发射质量为$m$，电荷量为$q(q>0)$的带电粒子，所有粒子的初速度大小均为$v$，粒子的发射速度的方向介于$y$轴和直线$y=\dfrac34x$之间，带电粒子经磁场偏转后均从$OP$边射出磁场，已知沿$y$轴正方向射出的粒子恰好通过$P$点。不计带电粒子的重力和粒子间的相互作用。

(1)求$OP$边界上有粒子射出的部分的长度$l$；

(2)若改变$OPMN$区域内的磁场分布范围，让磁场仅分布在$y$轴与满足某函数方程的边界之间，其他条件均不变，此时从粒子源$O$射出的粒子最终均从$MP$边界离开$OPMN$区域且通过$MP$边界时的速度方向相同，并且沿$y$轴正方向射出的粒子恰好通过$M$点，求

(i)粒子射出磁场时的方向；

(ii)磁场的边界方程。

%FIG p208_a 0.585 0.27 0.94 0.475
\notefig[0.42]{figures/p208_a.png}
\begin{solution}
(1)\ $r=\dfrac{a}{2}$\quad \unclear{由两个临界}\quad $l=\dfrac{\unclear{2}}{5}a$

(2)\ \red{$\theta=2\alpha$}

\red{$\sin\alpha=\dfrac{\sqrt5}{5}$}

\red{$\sin\theta=2\sin\alpha\cos\alpha=\dfrac45$}

\red{粒子在磁场中偏转角度为$53^\circ$，出磁场后速度方向与直线$y=\dfrac34x$平行}

(3)
{\color{red!75!black}
\[
\begin{cases}
r\cos\beta=x+r\cos\theta\\
r\sin\beta+y=r\sin\theta
\end{cases}
\]
\[
\left(\frac{x+r\cos\theta}{r}\right)^{2}+\left(\frac{r\sin\theta-y}{r}\right)^{2}=1\ \text{得}
\]
\[
\left(\frac{2x}{a}+\frac35\right)^{2}+\left(\frac45-\frac{2y}{a}\right)^{2}=1
\]
% 末行行尾的“1”被图旁线条遮住一部分，按上式补
}
\end{solution}
\end{problem}
%%%END
%%%BLOCK id=221-05 ch=11 sec=磁场·多解问题 kind=method alt=02,04 cont=none y=0.44-0.51
\textbf{5、}$r_{\text{轨}}$与几何长度关系不明 $\Leftrightarrow$ 多解、奇、偶

$f$的双向性
%%%END
%%%BLOCK id=225-09 ch=11 sec=复合场·带电体的直线运动 kind=method alt=02 cont=none y=0.70-0.90
%FIG p225_e 0.02 0.705 0.25 0.898
\notefig[0.3]{figures/p225_e.png}
\red{垂直于运动方向的合\unclear{外}力为0。} % 原稿：“合”与“力”之间有一处被涂的字迹

\red{正电\quad 电场向右。}
%%%END
%%%BLOCK id=225-10 ch=11 sec=复合场·摩擦与最大速度 kind=problem alt=03,09 cont=none y=0.78-1.00
\begin{problem}
%FIG p225_f 0.012 0.888 0.24 1.0
\notefig[0.3]{figures/p225_f.png}
求 $V_{\max}$ 和 $a_{\max}$ 时的 $V$
%FIG p225_g 0.285 0.885 0.475 0.995
\notefig[0.25]{figures/p225_g.png}
\begin{solution}
$a_{\max}$：$f=0$ % 原稿冒号处近似“÷”

$V_{\max}$：$a=0$

%FIG p225_h 0.603 0.806 0.93 0.876
\notefig[0.4]{figures/p225_h.png}
\noindent
\begin{minipage}[t]{0.48\linewidth}
开始 $V_{\text{小}}$
\begin{align*}
N+qVB-qE&=0\\
N&=qE-qVB\\
f&=\mu N\\
mg-f&=ma\\
a&=g-\frac{\mu qE}{m}+\frac{qVB}{m}\quad V\uparrow\ a\uparrow
\end{align*}
\end{minipage}\hfill
\begin{minipage}[t]{0.48\linewidth}
之后 $qVB>qE$
\begin{align*}
N+qE&=qVB\\
N&=qVB-qE\\
f&=\mu N\\
mg-f&=ma\\
a&=g+\frac{\mu qE}{m}-\frac{qVB}{m}\quad V\uparrow\ a\downarrow
\end{align*}
\end{minipage}
% CHECK: 两处 a 的表达式中第二项分子原稿均为 qVB（未见 μ）
\end{solution}
\end{problem}
%%%END
%%%BLOCK id=228-01 ch=11 sec=洛伦兹力分解·动量定理 kind=method alt=none cont=none y=0.03-0.17
\textbf{1.} 正则动量 \quad $qBx=m\Delta V_y$ \quad \red{2018、桥梁公式考} \quad \red{25T}
%FIG p228_a 0.09 0.075 0.28 0.165
\notefig[0.25]{figures/p228_a.png}

$qB_1l_1+qB_2l_2=m\Delta V_y=m(V_0\cos\alpha+V_0\sin\alpha)$
%%%END
%%%BLOCK id=228-02 ch=11 sec=复合场·多层电磁场 kind=problem alt=09 cont=none y=0.17-0.40
%FIG p228_b 0.03 0.175 0.235 0.265
\notefig[0.25]{figures/p228_b.png}

\begin{problem}
①求在第$n$层时出B的角度$\theta$（与水平方向）

②粒子恰可从第$n$层过，但比荷大一点不能穿过。解释
\begin{solution}
① $nqBd=m\Delta V_y$

$\Delta V_y=V\sin\theta_n$

$nEqd=\dfrac12 mV_n^{2}$ \qquad $V_n=\sqrt{\dfrac{2nEqd}{m}}$

$\Delta V_y=\sqrt{\dfrac{2nEqd}{m}}\,\sin\theta_n$

$\therefore\sin\theta_n=\sqrt{\dfrac{nqBd}{2Em}}$ % CHECK: 按推导根号内分子应为 nqB^2 d，原稿疑漏平方

② $\sin\theta\le1$ \qquad $\dfrac{q}{m}\uparrow$ \qquad $d\downarrow$

$d\downarrow$ 故出不去
\end{solution}
\end{problem}
%%%END
%%%BLOCK id=228-06 ch=11 sec=安培力·有效长度与张力 kind=note alt=02 cont=none y=0.73-0.85
%FIG p228_f 0.015 0.735 0.155 0.845
\notefig[0.25]{figures/p228_f.png}

$F_{\text{全}}=2BIL$

\underline{$F_{\text{半}}=\dfrac{F}{2}=BIL$} \struck{张力} \quad \red{上下张力合$=F=2BIL$} % 原稿 F半 式下有红色下划线
%%%END
%%%BLOCK id=239-02 ch=11 sec=带电粒子在磁场中的圆周运动·转筒 kind=note alt=none cont=none y=0.17-0.25
磁场转筒　$\omega_0=\omega_{\text{粒子}}$ % CHECK: 手写像 W₀=W，按角速度 ω 转录
%%%END
%%%BLOCK id=245-04 ch=11 sec=洛伦兹力·杆上带电小环 kind=problem alt=03 cont=none y=0.48-0.60
\begin{problem}
%FIG p245_f 0.0 0.49 0.20 0.605
\notefig[0.25]{figures/p245_f.png}
\begin{solution}
以$V$运动无弹力匀速直线运动

若$V_0>V$\quad 减速到$V_0$后匀速 % CHECK: 按物理应为“减速到V后匀速”，原稿写作 V₀

$a\downarrow$
\end{solution}
\end{problem}
%%%END
%%%BLOCK id=267-04 ch=11 sec=洛伦兹力·动量定理 kind=problem alt=06 cont=none y=0.48-0.60
\begin{problem}
%FIG p267_c 0.03 0.478 0.255 0.59
\notefig[0.3]{figures/p267_c.png}
\begin{solution}
\[ -BqV_xt-mgt=0-mV_0 \]
$V_{\text{末}}$ 用动能定理求
\[ t=\dfrac{mV_0-qB\dfrac{L}{2}}{mg} \]
\end{solution}
\end{problem}
%%%END
%%%BLOCK id=269-01 ch=11 sec=组合场·电场与磁场交替 kind=problem alt=09 cont=none y=0.02-0.67
% 贴纸印刷题；原稿上有圈画：“正”被圈出，“（不计重力）”“现只改变粒子自A点出射速度大小至v，”“粒子经过一段时间运动后可经过C点，”被划线；左侧页边另有难辨的铅笔涂写
\begin{problem}[13（贴纸印刷题）]
空间存在平面直角坐标系$xOy$，在$x<0$区域内有沿$x$轴正向的匀强电场，在$x>0$区域内有垂直平面向外的匀强磁场，在第二象限内有矩形OACD，$OA=\sqrt3h$，$OD=\unclear{2h}$，\unclear{一}个质量为$m$，电量为$q$的带正电粒子（\underline{不计重力}）从A点沿$y$轴正方向以某速度射入第二象限，经$t_0$时间后由D点进入磁场，又经一段时间射出磁场后又回到A点。\underline{现只改变粒子自A点出射速度大小至$v$，}\underline{粒子经过一段时间运动后可经过C点，}则（　　）

手写答案：AC

\begin{itemize}
\item[A.] 匀强电场的场强大小为$\dfrac{2\sqrt3\,mh}{qt_0^2}$ % 原稿：分式上有一道斜划线
\item[B.] 匀强磁场的磁感应强度大小为$\dfrac{\sqrt3\,m}{3qt_0}$ % 原稿：分式上有一道斜划线
\item[C.] 能使粒子以最短时间从A点运动至C点的初速度$v_0'=\dfrac{3h}{t_0}$ % 原稿：$\frac{3h}{t_0}$被大圈圈住
\end{itemize}

手写批注（选项旁/图旁）：
\begin{itemize}
\item \struck{$t_0=\dfrac{m}{qB}$} % 原稿：被圈出并被多道斜线划去
\item $qv_0Bt=mv_{\unclear{p}}$ % 原稿为“qV₀Bt=mV_ρ”样，V上有斜线，下标难辨
\item $\dfrac{2\sqrt3}{2}$ % 原稿：被斜线划去，位于选项C右侧
\item 页边左上铅笔涂写：\unclear{…}；\unclear{$\dfrac{\sqrt3}{2}$}
\end{itemize}

%FIG p269_a 0.62 0.158 0.93 0.352
\notefig[0.3]{figures/p269_a.png}

\begin{solution}
\[
\sqrt3h=\frac12at_0^2,\qquad a=\frac{qE}{m}
\]
\[
E=\frac{2\sqrt3\,mh}{qt_0^2}
\]
\[
R=\frac{2h}{\sin60^\circ}=\frac{mV}{qB}\qquad V=\frac{V_0}{\cos60^\circ}=\frac{4h}{t_0}\qquad\therefore B=\frac{\sqrt3\,m}{qt_0}
\]
\[
V_x=\frac{2\sqrt3h}{t_0}
\]
\[
\sin\theta=\frac{V_x}{V'}=\frac{2\sqrt3h}{V't_0}\qquad qV'B=\frac{mV'^2}{R}\qquad R'=\frac{mV'}{qB}\qquad R'\sin\theta=2h
\] % CHECK: qV'B=mV'^2/R 中的R疑应为R'
\[
y=2R'\sin\theta=4h\quad\text{（由对称性）}\quad\therefore y_1=3h\quad y_2=-h\quad\therefore V_0'=\frac{y_1}{t_0}=\frac{3h}{t_0}
\]
\end{solution}
\end{problem}
%%%END
%%%BLOCK id=269-02 ch=11 sec=正则动量 kind=note alt=none cont=none y=0.69-0.72
\textbf{正则动量}
%%%END
%%%BLOCK id=281-02 ch=11 sec=洛伦兹力·圆环套杆 kind=diagram alt=none cont=none y=0.19-0.34
%FIG p281_b 0.10 0.195 0.60 0.34 soft
\notefig[0.5]{figures/p281_b.png}
\unclear{有场}% 左侧示意图上方的淡铅笔小字；示意图中的点表示磁场方向，圆环套在杆上

图线旁标注：$qvB=mg$，$qvB>mg$，$qvB<mg$
%%%END
%%%BLOCK id=284-01 ch=11 sec=复合场·电场磁场交替区域 kind=problem alt=09,04 cont=none y=0.00-0.40
% 题干为粘贴的印刷题；手写对“E1”“E2”“匀强磁场”“正电”“初速度v0竖直向上抛出”“(-L,0)”画了方框/圈，对“区域II中做匀速圆周运动”“恰好通过坐标原点O回到区域I”“交替在区域I和区域II中运动”“第二次进入区域II时的位置坐标”“轨迹半径r2”画了下划线
% 右下角折页露出背面一小幅手绘图（斜线与一条末端带圈的曲线），属另一面内容，未转录
\begin{problem}[（20分）]
如图所示，在竖直平面内建立直角坐标系 $xOy$，$x$ 轴上方的区域I内存在水平向右的匀强电场 $E_1$，$x$ 轴下方的区域II内存在竖直向上的匀强电场 $E_2$ 和垂直于竖直平面向外的匀强磁场。将一个带正电的小球从 $x$ 轴上坐标为 $(-L,0)$ 的 $P$ 点以初速度 $v_0$ 竖直向上抛出，小球垂直通过 $y$ 轴进入第一象限，然后通过 $x$ 轴进入区域II中做匀速圆周运动，恰好通过坐标原点 $O$ 回到区域I，之后小球可以交替在区域I和区域II中运动。已知小球第三次进入区域II时的动能为其初动能的4倍，重力加速度未知。求：

（1）$E_1$ 和 $E_2$ 的大小之比；\quad （手写）\unclear{$\sqrt{3}:6$}% CHECK: 手写为“√3 ∶ 6”，6 的起笔很长，也可能是“√3 ∶ √6”

（2）小球第一次在区域II中运动的时间 $t_1$；\quad （手写）$\dfrac{2\pi L}{v_0}$

（3）小球第二次进入区域II时的位置坐标和第二次在区域II中做匀速圆周运动的轨迹半径 $r_2$

（手写）$r_2=\dfrac{\sqrt{21}}{2}L$

%FIG p284_a 0.645 0.015 1.0 0.235
\notefig[0.35]{figures/p284_a.png}
\end{problem}
%%%END
%%%BLOCK id=285-01 ch=11 sec=带电粒子在电磁场中的运动·三类场 kind=summary alt=09 cont=none y=0.00-0.07
% 页面上缘被裁，首行只露出下半截；以下三行由左侧一条竖线（大括号）括起，竖线左边为标题
\unclear{图带规律}
\begin{itemize}
\item 电偏转\ \unclear{…}
\item 复合场\quad $v=\dfrac{E}{B}$
\item 磁偏转\quad $r=\dfrac{mv}{qB}$
\end{itemize}
%%%END
%%%BLOCK id=285-03 ch=11 sec=带电粒子在磁场中的圆周运动·圆心确定 kind=method alt=none cont=none y=0.08-0.12
初速度垂线\quad 弦长中垂线\quad $V_{\text{初}}$、$V_{\text{末}}$ 的\unclear{角}平分线\quad \fbox{交于圆心}
%%%END
%%%BLOCK id=289-04 ch=11 sec=洛伦兹力分解·动量定理 kind=problem alt=06 cont=none y=0.38-0.60
%FIG p289_d 0.055 0.38 0.25 0.495
\notefig[0.28]{figures/p289_d.png}

\begin{gather*}
-mgt-qBV_xt=0-mV_0\\
mgt+\frac12 Bql=\text{\struck{$0-$}}\,mV_0\qquad t=\frac{mV_0-\frac12 qBL}{mg}\\
-mgd=\frac12 m\left(V_{\text{末}}^2-V_0^2\right)
\end{gather*}
% CHECK: 第三式 V 的下标手写像“末”；第二式中 l 与 t 式中 L 应为同一量（图中标 L/2）；第二式等号后被划去的是“0-”
%%%END
%%%BLOCK id=289-05 ch=11 sec=洛伦兹力·环套杆摩擦做功 kind=problem alt=06 cont=none y=0.61-0.78
%FIG p289_e 0.10 0.635 0.27 0.72
\notefig[0.28]{figures/p289_e.png}

%FIG p289_f 0.295 0.635 0.86 0.78
\notefig[0.55]{figures/p289_f.png}

\begin{itemize}
\item $mg=qVB$\quad $W_f=0$
\item $mg\le qVB$\quad $W_f=\text{\unclear{…}}-\dfrac12\left(mV_0^2-V_{\text{末}}^2\right)$
\item $mg>qVB$\quad $W_f=\dfrac12 mV_0^2$
\end{itemize}
% CHECK: 第二种情形 W_f 等号后有一个无法辨认的涂改小记号（像被划去的 E/½）；括号内第二项漏写 m；V 的下标手写像“末”；“mg≤qVB”也可能是 mg<qVB
%%%END
%%%BLOCK id=290-01 ch=11 sec=带电粒子在磁场中的圆周运动·打板平均作用力 kind=problem alt=07 cont=none y=0.01-0.27
%FIG p290_a 0.065 0.015 0.29 0.12
\notefig[0.25]{figures/p290_a.png}

$V_0\sim\in(0,\sqrt{3}V_0)$% CHECK: 首行第一个符号 V 的下标及其后的“~”不确定（也可能是 $V_y\in(0,\sqrt3V_0)$）

每秒 $N_0$ 个粒子射出。\quad $U=\dfrac{mV_0^2}{2q}$\ $\cdot$\ $d=\dfrac{mV_0}{qB_0}$

$V_y\in(V_0,2V_0)$\quad $N=\dfrac{3d-2d}{\underset{\uparrow\,\text{\unclear{能到}}}{4d}-\underset{\uparrow\,\text{\unclear{能到}}}{2d}}N_0=\dfrac12N_0$
% 分母 4d、2d 下方各有一个向上的箭头指向该项，箭头下各写“能到”（字形不清）

平均作用力\quad $r_3=\dfrac32d$\quad 打在中间% CHECK: “r_3”手写像 γ_3（也可能是 y_3），按上下文读作 r_3

\[ \bar V=\frac{V_1+V_3}{2} \]
\[ Ft=Nm\bar V=\frac58N_0mV_0 \]
%%%END
%%%BLOCK id=290-02 ch=11 sec=圆形磁场·磁聚焦与偏转电场 kind=problem alt=09,07 cont=none y=0.28-0.54
% 剪贴的印刷题；原稿中被手工框出的词句已用 \fbox 标出（“板间距离”与“也为 r”、“单位时间”与“内发射的粒子数为 n”在原稿中分处两行）；题图上另有手绘的粒子轨迹
\begin{problem}
（18分）如图所示，真空中有一半径为 $r$ 的圆形磁场区域，与 $x$ 轴相切于坐标原点 $O$，\fbox{磁感应强度大小为 $B$}，方向\fbox{垂直于纸面向外}，在圆形磁场区域的右侧有两水平放置的正对平行金属板 $M$、$N$，\fbox{板间距离也为 $r$}，\fbox{板长为 $L$}，两板的中间线 $O_2O_3$ 的反向延长线恰好过磁场圆的圆心 $O_1$。在 $O$ 点处有一粒子源，能向磁场中各个方向（纸面内）源源不断发射\fbox{速率相同}、\fbox{质量为 $m$、比荷为 $k$} 且带\fbox{正电}的粒子，\fbox{单位时间内发射的粒子数为 $n$}，沿 $y$ 轴正方向射入磁场的粒子，恰能沿直线 $O_2O_3$ 射入平行金属板间。不计粒子重力、空气阻力及粒子间的相互作用。

(1) 求粒子射入磁场的速率 $v_0$。

(2) 若已知两平行金属板间电场强度 $E=\dfrac{kB^2r^3}{L^2}$，在平行板的右端适当位置固定一平行于 $y$ 轴方向的探测板（图中未画出），使从 $M$、$N$ 板右端射出平行板间的粒子全部打在探测板上，则探测板板长至少为多少？

(3) 在满足第 (2) 问的前提下，若打到探测板上的粒子被全部吸收，求单位时间粒子束对探测板的平均作用力的大小。

%FIG p290_b 0.312 0.398 0.595 0.5
\notefig[0.3]{figures/p290_b.png}
\begin{solution}
(1)\quad $V_0=kBr$

(2)\quad $\left\{\begin{aligned} y&=\tfrac12at^2\\ t&=\tfrac{L}{V_0}\qquad y=\tfrac12r\\ a&=\tfrac{qE}{m}\end{aligned}\right.$

(3)
\begin{align*}
V_x&=kBr & \tan\theta&=\frac{V_y}{V_x}=\frac{2\cdot\frac{r}{2}}{L}\\
V_y&=\frac{r}{L}\,kBr
\end{align*}
\[ F=\frac13nm\,\frac{\sqrt{V_x^2+V_y^2}}{2}=\frac{kBr}{6L}\sqrt{r^2+L^2} \]
% CHECK: 根号内手写像 r³L²，按上下文读作 r²+L²；F 的第一个根号下“2”的位置（√(…)/2 还是 √(…/2)）按结果 1/(6L) 读作 √(…)/2；等号右端未写 nm
\end{solution}
\end{problem}
%%%END
%%%BLOCK id=293-01 ch=11 sec=有界磁场·动态圆 kind=method alt=none cont=none y=0.03-0.20
\textbf{旋转圆}

%FIG p293_a 0.04 0.052 0.245 0.195
\notefig[0.3]{figures/p293_a.png}

$v=\dfrac{qBd}{m}$\quad 定速不定向

先交后切\quad 弦最短时 $t$ 最短(半径)。
%%%END
%%%BLOCK id=293-02 ch=11 sec=有界磁场·动态圆 kind=method alt=none cont=none y=0.16-0.33
\textbf{缩放圆}\quad 定向不定速

%FIG p293_b 0.056 0.187 0.345 0.335
\notefig[0.35]{figures/p293_b.png}

\begin{center}
\begin{tabular}{@{}l@{\hspace{2.5em}}l@{}}
打在左 & 上 \\
$[0,\frac{1}{3}L]$ & $(\frac{1}{3}L,(4-2\sqrt{3})L]$ \\[3pt]
$((4-2\sqrt{3})L,L]$ & $(L,+\infty)$ \\
右 & 下
\end{tabular}
\end{center}
%%%END
%%%BLOCK id=293-03 ch=11 sec=有界磁场·动态圆 kind=method alt=none cont=none y=0.33-0.54
\textbf{汇聚圆}\quad $R=d=r$

%FIG p293_c 0.082 0.332 0.40 0.54
\notefig[0.35]{figures/p293_c.png}

图旁标注：菱形
%%%END
%%%BLOCK id=293-04 ch=11 sec=有界磁场·动态圆 kind=method alt=none cont=none y=0.54-0.63
\textbf{平移圆}\quad 定速定向不定心

%FIG p293_d 0.11 0.562 0.275 0.63
\notefig[0.25]{figures/p293_d.png}
%%%END
%%%BLOCK id=295-04 ch=11 sec=有界磁场·圆形边界 kind=problem alt=none cont=none y=0.55-0.83
% 以下题干为剪贴的印刷题（题号被裁掉，仅剩一个“.”），其上有手写划线/圈画/解答
\begin{problem}
如图所示，在以 $R_0$ 为半径，$O$ 为圆心的圆形区域内存在磁场，直径 $MN$ 左侧区域存在\underline{一匀强磁场}，方向垂直于纸面向外，磁感应强度大小为 $B_1$；$MN$ 右侧区域也存在一匀强磁场，方向垂直于纸面向里，磁感应强度大小为 $B_2$，有一质量为 $m$，电荷量\underline{为$+q$ 的}带电粒子（不计重力）沿垂直于 $MN$ 的方向从 $P$ 点射入磁场，通过磁场区域后自 $Q$ 点离开磁场，离开磁场时其运动方向仍垂直于 $MN$。已知 $OP$ 与 $MN$ 的夹角为 $\theta_1$，$OQ$ 与 $MN$ 的夹角为 $\theta_2$，粒子在左侧区域磁场中的运动时间为 $t_1$，粒子在右侧区域磁场中的运动时间为 $t_2$，则下列说法\fbox{正确的}是（\quad）

\begin{tabular}{@{}l@{\hspace{3em}}l@{}}
A. $\dfrac{B_2}{B_1}=\dfrac{\sin\theta_1}{\sin\theta_2}$ & B. $\dfrac{B_2}{B_1}=\dfrac{\sin\theta_2}{\sin\theta_1}$ \\[8pt]
C. $\dfrac{t_1}{t_2}=\dfrac{\sin\theta_2}{\sin\theta_1}$ & D. $\dfrac{t_1}{t_2}=\dfrac{\sin\theta_1}{\sin\theta_2}$
\end{tabular}
% 手写：选项 A 的分子处有一勾状弧线，选项 D 的 t1/t2 上有一道斜线（疑为勾选）

%FIG p295_e 0.76 0.653 0.94 0.812
\notefig[0.3]{figures/p295_e.png}

\begin{solution}
手写答案：（A）D\quad $r_1\propto\sin\theta_1$，

作垂线 % 其下有一条弧线伸向右侧

$\sin\theta_2=\dfrac{d_2}{R}$\quad $r_2=\dfrac{d_2}{\sin\alpha}$

$\sin\theta_1=\dfrac{d_1}{R}$\quad $r_1=\dfrac{d_1}{\sin\alpha}$
\end{solution}
\end{problem}
%%%END
%%%BLOCK id=299-01 ch=11 sec=摆线运动 kind=diagram alt=09 cont=none y=0.04-0.22
% 笔记中速度多写作大写 V，此处统一转为小写 v
\textbf{初速为0}

%FIG p299_a 0.085 0.04 0.60 0.197
\notefig[0.5]{figures/p299_a.png}

$S=vT$.

\[ qE=qv_0B \]
%%%END
%%%BLOCK id=299-02 ch=11 sec=摆线运动 kind=diagram alt=09 cont=none y=0.24-0.45
\textbf{初速水平}

$v_0=v_2-v_1$

%FIG p299_b 0.08 0.24 0.66 0.41
\notefig[0.55]{figures/p299_b.png}

\[ qEh=\frac12 mv^2-\frac12 mv_0^2 \quad\text{或}\quad h=R(1-\sin\theta)\qquad R=\frac{mv}{qB} \]
%%%END
%%%BLOCK id=299-03 ch=11 sec=摆线运动 kind=diagram alt=09 cont=none y=0.455-0.64
\textbf{初速向上}

%FIG p299_c 0.09 0.455 0.56 0.64
\notefig[0.45]{figures/p299_c.png}
%%%END
%%%BLOCK id=299-04 ch=11 sec=摆线运动 kind=derivation alt=07 cont=none y=0.48-0.58
% 本块位于“初速向上”图的右侧
\red{常规方法 \unclear{正则动量加位移微积理}} % CHECK: 红字潦草，仅“常规方法”较有把握，其余字迹难辨
\[ \sum\Delta t\,qB\,\Delta v_x=\sum m a_y\cdot\Delta t \] % CHECK: 手写为“Σ△t qB△Vx=Σ m a_y·△t”，Vx 前的“△”也可能是圆点；物理上应为 Σ qBv_xΔt=Σ m a_yΔt
\[ qBx=m\Delta v_y \]
%%%END
%%%BLOCK id=299-05 ch=11 sec=有界磁场·动态圆 kind=summary alt=none cont=next y=0.67-1.00
% 红笔曲线从“弦切角”处起笔，向右延伸出页面右缘，接到下一页(p300)左缘的红色长曲线（指向 t_max 公式）
\textbf{临界问题}
\begin{itemize}
  \item 轨迹变化 $\to$ 动态圆 $\to$
  \begin{itemize}
    \item 直径作弦
    \item 相切。
  \end{itemize}
  \item 区域
  \begin{itemize}
    \item 直线边界\quad $L=2r\sin\theta$
    \item 直角边界
    \item 圆形边界
    \begin{itemize}
      \item $R=r$\quad 磁聚焦、/发聚
      \item $R>r$ \\ $R<r$ \quad 几个直角三角形\quad \red{圆心连线、公切线}
    \end{itemize}
  \end{itemize}
  \item \unclear{$t$}的范围 % CHECK: 手写首字形似“七”，按上下文读作 t
  \begin{itemize}
    \item 平移圆\red{/}旋转圆 $\to$ \red{\fbox{弦长}} % 红笔圈出“弦长”；“平移圆”与“旋转圆”之间有红色短竖线
    \item 放缩圆 $\to$ \red{\underline{弦切角}} % 红色下划线，并由此延伸出长红曲线
  \end{itemize}
\end{itemize}
%%%END
%%%BLOCK id=300-01 ch=11 sec=有界磁场·动态圆 kind=diagram alt=none cont=prev y=0.00-0.50
% 承接上一页(p299)“弦切角”处延伸过来的红色长曲线：自本页左缘 y≈0.85 处入页，向右上弯至 t_max 公式下方，箭头指向 t_max；页面下半部其余为空白横线（背面透字已忽略）
% 图为铅笔画的一族过同一点的圆，加红笔坐标轴、斜线、虚线、箭头和标注（均在图里）
%FIG p300_a 0.05 0.0 0.95 0.50
\notefig[0.85]{figures/p300_a.png}

\noindent\red{弦最长\unclear{为$\infty$}}\par
\noindent\red{直径当弦长}\par
\noindent\red{$\theta=2\alpha$\quad $t_{\max}=\dfrac{2\alpha m}{qB}$\quad 圆心角$=2\times$弦切角}
%%%END
